\section{General Duality Framework}
\label{sec:general_framework}

\paragraph{Setup from the full Hamiltonian.}
Take a subsystem \(Q\) and environment \(X\) with
\begin{equation}
  H_{\mathrm{tot}} = H_Q + H_X + H_{\mathrm{int}}, \qquad
  H_{\mathrm{int}} = g\,f\otimes B.
\end{equation}
Here \(g\) is a dimensionless coupling amplitude, \(f\) is a system operator, and \(B\) is a bath
operator (with bath mean absorbed so \(\Tr_X(B e^{-\beta H_X})=0\)). Define
\begin{equation}
  Z_X(\beta):=\Tr_X e^{-\beta H_X}, \qquad
  Z_Q(\beta):=\Tr_Q e^{-\beta H_Q},
\end{equation}
\begin{equation}
  \rho_Q^{\mathrm{th}}(\beta):= Z_Q(\beta)^{-1}e^{-\beta H_Q}.
\end{equation}
The unnormalized reduced equilibrium operator is
\begin{equation}
  \bar{\rho}_Q(\beta):=\Tr_X e^{-\beta H_{\mathrm{tot}}},
\end{equation}
with normalized state \(\rho_Q=\bar{\rho}_Q/\Tr_Q\bar{\rho}_Q\).

The Hamiltonian of mean force is defined exactly by
\begin{equation}
  e^{-\beta H_{\mathrm{MF}}(\beta)}
  := \frac{\bar{\rho}_Q(\beta)}{Z_X(\beta)}.
\end{equation}
This identity is purely definitional and non-perturbative
\cite{gelinThermodynamicsSubensembleCanonical2009a,hiltHamiltonianMeanForce2011,talknerColloquiumStatisticalMechanics2020,rivasStrongCouplingThermodynamics2020}.

\begin{definition}[Influence operator]
\label{def:influence_operator}
Define \(S_\beta\) by
\begin{equation}
  e^{S_\beta}
  := e^{+\beta H_Q/2}\,\frac{\bar{\rho}_Q(\beta)}{Z_X(\beta)}\,e^{+\beta H_Q/2}.
\end{equation}
Equivalently,
\begin{equation}
  \frac{\bar{\rho}_Q(\beta)}{Z_X(\beta)}
  = e^{-\beta H_Q/2}\,e^{S_\beta}\,e^{-\beta H_Q/2}.
\end{equation}
Hence
\begin{equation}
  \rho_Q(\beta)
  = \frac{Z_Q(\beta)}{Z_{\mathrm{MF}}(\beta)}
  \big(\rho_Q^{\mathrm{th}}(\beta)\big)^{1/2}
  e^{S_\beta}
  \big(\rho_Q^{\mathrm{th}}(\beta)\big)^{1/2},
\end{equation}
where \(Z_{\mathrm{MF}}(\beta):=\Tr_Q e^{-\beta H_{\mathrm{MF}}(\beta)}\).
\end{definition}

This thermal-plus-influence partition is the structural backbone used below: bare thermal factors
carry the reference Gibbs geometry, while \(S_\beta\) carries the interaction dressing. The duality
results concern representative maps acting on this dressed family. It is the operator-level
equilibrium counterpart of influence-functional dressing
\cite{feynmanTheoryGeneralQuantum1963a,mccaulMeanForceHamiltoniansInfluence2026}.

\begin{definition}[Representative family]
\label{def:rep_family}
Let \(\mathcal{P}\) be a parameter manifold and \(\mathcal{O}\) a set of observables. A
representative family is a map
\begin{equation}
  \mathcal{R}: \mathcal{P}\to \mathcal{O}, \qquad
  p \mapsto \{O_k(p)\}_{k\in K},
\end{equation}
where each \(O_k\) is evaluated from one fixed analytic functional form, while \(p\) is allowed
to change by discrete representative transformations.
\end{definition}

\begin{definition}[Crossover coordinate]
Assume there exists a positive scalar control coordinate \(\chi\) and its logarithmic coordinate
\(y\),
\begin{equation}
  \chi = \chi(\beta,\lambda,\vartheta) > 0,
\end{equation}
\begin{equation}
  y := \log \chi.
\end{equation}
In the single-channel setting used below,
\begin{equation}
  \chi(\beta,g,\vartheta)=g^2\chi_0(\beta,\vartheta), \qquad \chi_0>0.
\end{equation}
\end{definition}

\begin{definition}[Admissible observable class]
\label{def:admissible_class}
An observable \(Q\) is admissible if it can be written as
\begin{equation}
  Q(\beta,g,\vartheta)=F\!\bigl(\beta,\vartheta;\phi_1(\chi),\ldots,\phi_n(\chi),u\bigr),
\end{equation}
where \(F\) is smooth in its arguments, \(\phi_j\) have asymptotic expansions for
\(\chi\to 0\) and \(\chi\to\infty\), and \(u\) is an oriented scalar (possibly absent) that is odd
under orientation-gauge flip.
\end{definition}

\begin{proposition}[Representative-map admissibility]
\label{prop:admissibility}
If \(Q\) belongs to the admissible class in Definition~\ref{def:admissible_class}, then discrete
maps acting on representative labels \((\chi,u)\) may exchange the computational asymptotic regime
of \(Q\) while preserving the same analytic functional form \(F\).
\end{proposition}
\begin{proof}
The map acts only on the arguments of \(F\), not on \(F\) itself. Hence transformed evaluations are
obtained by argument substitution and remain within the same closed-form family.
\end{proof}

\paragraph{Theory/data separation.}
The next three sections use only the structure above. Model closure and bath-specific formulas enter
only in the spin-boson example section.
